\scanpage{121}
% archive sheet number: 120
\diagram{\begin{tikzcd}[ampersand replacement=\&]
 B\arrow[d,no head]\&k^B\arrow[d,"t=\mathrm{Tr}"]\arrow[r,leftrightarrow,"\unsure{\varphi|\rho}=(\varphi_b)_{b\in B}"]\&E=t^{-1}(\{1\})\\
 I\&k^I\&
\end{tikzcd}}
\[
 H_b=\varphi_b^{-1}(\{1\})
\]
\diagram{\begin{tikzcd}[ampersand replacement=\&]
 B'\subset B\arrow[d]\&\deleted{\unsure{A}}F_{B'}=\displaystyle\bigcap_{b\in B'}H_b\\
 I\&
\end{tikzcd}}
---
\[
 \Phi=\{F_{B'}\}_{B'\in\Gamma_p(B/I)}\cup\{\varnothing\}
\]
\noindent\rule{\linewidth}{0.4pt}

\diagram{\begin{tikzcd}[ampersand replacement=\&]
 \deleted{\unsure{E}}N\arrow[r,hook]\&k^B\arrow[r]\&k
\end{tikzcd}}
$N=$ sous-module (facteur-direct de rang $n-1$) engendré par les $(b+\sigma b)-(b'+\sigma b')$ $b,b'\in B$
\diagram{\begin{tikzcd}[ampersand replacement=\&]
 0\arrow[r]\&\underset{n}{W}\arrow[r]\&V\arrow[r]\arrow[d,equal]\&k\arrow[r]\&0\\
 \&\&k^B/N\&\&
\end{tikzcd}}

\uline{Caractérisation du cocube}

$\Phi\subset E$ espace affine de dim $n$
\inserted{ens de sous-esp. \unsure{lin.-affines}}

\begin{enumerate}[label=\alph*)]
\item $\forall s\in\Phi_0$ (él. de $\Phi$ qui est un sous-espace de $E$ de dim 0, i.e. réduit à un point) $\exists!s'\in\Phi_0$, \inserted{$s'\ne s$,} tel que la droite joignant $s,s'$ soit $\notin\Phi$ ($s\mapsto s'$ est alors une involution ss pt fixe en $\Phi_0$) \inserted{not $\sigma$}
\item Pour tte partie $B'\subset B\overset{(\mathrm{dfn})}=\Phi_0$, telle que $B'\in\Gamma_p(B/I)$ (où $I=B/\sigma$), le sous-espace \inserted{lin.} affine engendré \unsure{pr|par} $B'$ est $\in\Phi$ --- \inserted{et de dim égale à $(\mathrm{card}\,B')-1$ (i.e. $B'$ est aff\textsuperscript{t} libre)} et on trouve ainsi tous les $\unsure{\in|f\in}\Phi$, \deleted{sous-espaces affines}, à l'exception de $E$ (l'élément maximal), qui est $\in\Phi$
\item Pour toute $B'\in\Gamma(B/I)$ et $b\in B\setminus B'$, $B'\cup\{b\}$ est aff\textsuperscript{t} libre et engendre aff\textsuperscript{t} $E$
\item Si $b,b'\in B$, on a $b+\sigma b=b'+\sigma b'$.
\end{enumerate}
